%%%%%%%%%%%%%%%%%%%%%%%%%%%%%%%%%%%%%%%%%%%%%%%%%%%%%%%
% Arquivo de entrada de dados da parte pré-textual
%%%%%%%%%%%%%%%%%%%%%%%%%%%%%%%%%%%%%%%%%%%%%%%%%%%%%%%

% Tipo do trabalho
%-----------------------------------------------------------------------------
\tipotrabalho{Trabalho de Conclusão de Curso}

% Elementos que compõem o documento
%-----------------------------------------------------------------------------
 \providecommand{\terCapa}{Sim}
 \providecommand{\terFolhaRosto}{Sim}
 \providecommand{\terTermoAprovacao}{Sim}
 \providecommand{\terDedicatoria}{Nao}
 \providecommand{\terFichaCatalografica}{Nao}
 \providecommand{\terEpigrafe}{Nao}
 \providecommand{\terAgradecimentos}{Sim}
 \providecommand{\terErrata}{Nao}
 \providecommand{\terListaFiguras}{Sim}
 \providecommand{\terListaQuadros}{Nao}
 % Sem quadros no documento: todo elemento tabular e apresentado
 % como tabela, e a Lista de Tabelas os reune em serie unica.
 \providecommand{\terListaTabelas}{Sim}
 \providecommand{\terSiglasAbrev}{Sim}
 \providecommand{\terSimbolos}{Sim}
 \providecommand{\terResumos}{Sim}
 \providecommand{\terSumario}{Sim}
 \providecommand{\terAnexo}{Nao}
 \providecommand{\terApendice}{Sim}
 \providecommand{\terIndiceR}{Nao}
%-----------------------------------------------------------------------------

%% TODO: dado ausente — a ficha catalográfica é emitida pelo Sistema de
%% Bibliotecas da UFPR. Enquanto não for emitida, \terFichaCatalografica
%% permanece em "Nao"; ao recebê-la, gravar o arquivo como
%% metadados/FichaC.pdf e alterar a marca para "Sim".

% Identificação
%-----------------------------------------------------------------------------
\titulo{Dimensionamento automatizado de cobertura para estacionamento em
madeira lamelada colada segundo o Eurocódigo 5}

\autor{Gabriel Goes Stocki}
\local{Curitiba}
\data{2026} % Apenas ano, 4 dígitos

% Orientação
%-----------------------------------------------------------------------------
% Preenche-se apenas o comando correspondente. Havendo os dois, prevalece o
% feminino; por isso \orientadora permanece vazio.
\orientador{Prof. Dr. Elvidio Gavassoni Neto}
\orientadora{}

% Titulação informada pelo autor (2026-11-14): Doutor e Mestre em Engenharia
% Estrutural (PUC-Rio); Engenheiro de Terminais e Dutos (Universidade Petrobras);
% Engenheiro Civil (UFES); Agrimensor (ETFES). Na folha de rosto e no termo de
% aprovação entra só o título acadêmico mais alto, abreviado antes do nome:
% "Prof. Dr." (doutorado). Os demais não se indicam nos pré-textuais.

% Não há coorientação neste trabalho.
\coorientador{}
\coorientadora{}
\scoorientador{}
\scoorientadora{}

% ----------------------------------------------------------
\addbibresource{referencias.bib}

% Instituição
%-----------------------------------------------------------------------------
\instituicao{%
Universidade Federal do Paraná}

\def \ImprimirSetor{Setor de Tecnologia}

\def \ImprimirProgramaPos{}

\def \ImprimirCurso{Pós-Graduação \textit{Lato Sensu} em Estruturas}

\preambulo{%
Trabalho de Conclusão de Curso apresentado como requisito parcial à obtenção do
título de Especialista em Estruturas, Pós-Graduação \textit{Lato Sensu} em
Estruturas, Setor de Tecnologia, Universidade Federal do Paraná}

%% Denominação do curso e título conferido informados pelo autor (issue 67).

%-----------------------------------------------------------------------------
\newcommand{\imprimirCurso}{}

\newcommand{\imprimirDataDefesa}{%
\sugest{11/12/2026}}

%% TODO: dado ausente — data da defesa.

\newcommand{\imprimircdu}{%
\sugest{[CDU]}}

%% TODO: dado ausente — a CDU é atribuída pelo Sistema de Bibliotecas da UFPR
%% no momento da emissão da ficha catalográfica.

% ----------------------------------------------------------
\newcommand{\imprimirerrata}{}

% Comandos de dados — Data da apresentação
\providecommand{\imprimirdataapresentacaoRotulo}{}
\providecommand{\imprimirdataapresentacao}{}
\newcommand{\dataapresentacao}[2][\dataapresentacaoname]{\renewcommand{\dataapresentacao}{#2}}

% Comandos de dados — Nome do Curso
\providecommand{\imprimirnomedocursoRotulo}{}
\providecommand{\imprimirnomedocurso}{}
\newcommand{\nomedocurso}[2][\nomedocursoname]
  {\renewcommand{\imprimirnomedocursoRotulo}{#1}
\renewcommand{\imprimirnomedocurso}{#2}}

% Banca examinadora — termo de aprovação
%-----------------------------------------------------------------------------
%% TODO: dado ausente — substituir os marcadores pelos nomes e pelas
%% instituições dos membros da banca examinadora.
\newcommand{\AssinaAprovacao}{

% A assinatura do orientador é emitida automaticamente por \insereAprovacao;
% aqui entram apenas os demais membros da banca.
\assinatura{\sugest{[Nome do segundo membro]} \\ \sugest{[Instituição]}}
   \assinatura{\sugest{[Nome do terceiro membro]} \\ \sugest{[Instituição]}}

   \begin{center}
    \vspace*{0.5cm}
    \imprimirlocal, \imprimirDataDefesa.
    \vspace*{1cm}
  \end{center}
  }

% Epígrafe
%-----------------------------------------------------------------------------
\newcommand{\EpigrafeTexto}{}

% Resumo e palavras-chave
%-----------------------------------------------------------------------------
\newcommand{\ResumoTexto}{%
Este trabalho apresenta o dimensionamento e a verificação estrutural de uma
cobertura para estacionamento em madeira lamelada colada, aberta nas quatro
faces, com 40,5 m por 46,0 m em planta, composta por nove pórticos de duas vigas de duplo-bisel de
20,25 m de vão apoiadas sobre três pilares. O dimensionamento segue a
EN 1995-1-1, complementada pela EN 1990 e pela EN 14080, com as forças devidas
ao vento determinadas pela ABNT NBR 6123. A cadeia de \nVerifExt{} verificações foi
implementada em programa de computador submetido a suíte de \nTestes{} testes
automatizados e conferido por casos de referência independentes, o que
permitiu avaliar 3.213 geometrias candidatas pelo mesmo conjunto de verificações
que produz o memorial final, com o mesmo código na fase de concepção e na de
verificação. A análise da rigidez relativa
entre a viga e os seus apoios demonstrou que o empuxo de arco não se
desenvolve, o que elevou o momento de projeto em 70,4\% em relação ao modelo de
arco e dispensou o tirante de base. Com vigas e pilares rotulados, a força
transversal só chega aos frontões por um diafragma com alma em todas as baías.
Na solução adotada, com os critérios de margem e os detalhes escolhidos, a maior
taxa da viga é de 0,58, no apoio sobre a sela central, e as maiores utilizações
estão nas peças de aço das ligações, nos contraventamentos e nas lâminas dos
pilares montantes, com até 0,91, e no critério de engajamento dos pinos da
ligação entre viga e pilar, com \etaMaxGlobal. A parcela de atrito do vento
paralelo à cumeeira, prescrita pela norma brasileira, mostrou-se determinante
para todo o sistema de estabilização longitudinal. Apoiar a viga por contato
direto sobre o pilar, deixando aos pinos apenas o arrancamento, reduziu o
fendilhamento nas ligações entre viga e pilar a uma taxa de 0,67, sem necessidade
de reforço. As ligações, a ancoragem e os demais elementos foram detalhados em
16 folhas de desenho cotadas e em escala, geradas pelo mesmo programa, que
fornece também os esforços destinados ao projeto da fundação. A estimativa
preliminar de custo da estrutura resultou em R\$ \custoTotalMi{} milhões, ou
R\$ \custoPorMetro{} por metro quadrado coberto, dos quais \custoPctMadeira~\%
correspondem à madeira. Conclui-se que, na solução estudada, os sistemas de
estabilização e os detalhes das ligações condicionaram o projeto tanto quanto a
resistência da viga, e que a automação da cadeia de verificação permite
conhecer a posição da solução adotada entre as geometrias aprovadas na grade
examinada.
}

\newcommand{\PalavraschaveTexto}{%
Madeira lamelada colada; Eurocódigo 5; dimensionamento estrutural; varredura
paramétrica; contraventamento.}

% Abstract and keywords
%-----------------------------------------------------------------------------
\newcommand{\AbstractTexto}{%
This work presents the structural design and verification of a glued laminated
timber parking canopy, open on all four sides, measuring 40.5 m by 46.0 m in plan
and comprising nine frames of two double-tapered beams spanning 20.25 m each,
supported on three columns. The design follows EN 1995-1-1, complemented by
EN 1990 and EN 14080, with wind actions determined according to the Brazilian
standard ABNT NBR 6123. The chain of \nVerif{} verifications was implemented in a
computer program subjected to a suite of \nTestes{} automated tests and checked
against independent reference cases, which made it possible to assess 3,213
candidate geometries using the same set of verifications that produces the
final design report, with the same code in the conceptual and the verification
stages. Analysis of the
relative stiffness between the beam and its supports showed that arch thrust
does not develop, which raised the design bending moment by 70.4\% with respect
to the arch model and rendered the base tie unnecessary. With pinned beams and
columns, the transverse force reaches the gable ends only through a roof
diaphragm with bracing in every bay. In the adopted solution, under the chosen
margin criteria and details, the highest beam utilisation is 0.58, in bearing on
the central saddle, and the highest utilisations are in the steel parts of the
connections, in the bracing systems and in the leaves of the bracing columns, at
up to 0.91, and in the dowel engagement criterion of the beam-to-column
connection, at 0.95. The wind friction component parallel to the ridge,
prescribed by the Brazilian standard, proved decisive for the entire
longitudinal stabilisation system. Supporting the beam by direct bearing on the
column, leaving only the uplift to the dowels, reduced splitting at the
beam-to-column connections to a utilisation ratio of 0.67, without
reinforcement. The connections, the anchorage and the remaining members were
detailed in 16 dimensioned, scaled drawing sheets produced by the same program,
which also delivers the design forces for the foundation design. The preliminary
cost estimate of the structure is BRL \custoTotalMiEn{} million, or
BRL \custoPorMetroEn{} per square metre of covered area, of which
\custoPctMadeira\% corresponds to the timber. It is concluded that, in the
solution studied, the stabilisation systems and the connection details
conditioned the design as much as the beam resistance, and that automating the
verification chain makes it possible to locate the adopted solution among the
geometries approved in the examined grid.
}

\newcommand{\KeywordsTexto}{%
Glued laminated timber; Eurocode 5; structural design; parametric sweep;
bracing.}

% ----------------------------------------------------------
\newcommand{\Resume}{}
\newcommand{\Motscles}{}
\newcommand{\Resumen}{}
\newcommand{\Palabrasclave}{}

% Agradecimentos
%-----------------------------------------------------------------------------
\newcommand{\AgradecimentosTexto}{%
Agradeço a Deus por me conceder saúde, à minha família por sempre acreditar em mim, à minha futura esposa pelo apoio incondicional e ao meu orientador por todos os ensinamentos ao longo dessa caminhada.

}

% Dedicatória
%-----------------------------------------------------------------------------
\newcommand{\DedicatoriaTexto}{}
